\section{The selfdual higher symplectic gate over three elements}
\label{sec:higher-symplectic-three}

\paragraph{Status and limitation.}
This new proof has passed complete mathematical reviews by /root and
/root/imprimitive. Both read the entire author checker and its output;
the peer also supplied a separately authored exact arithmetic certificate.
It uses the complete
opposite-radical and fixed-subspace carrier proofs D-HOS-07--08.
The normalizer bound below applies to every ordinary irreducible
representation. Its source consequence requires selfduality of the
irreducible core representations over $\mathbf F_3$. No reality of
all symplectic characters over this field is assumed. Nonselfdual
core factors remain outside the final obstruction.

\begin{lemma}[D-HOS-11: full normalizers over $\mathbf F_3$]
\label{D-HOS-11}
Let $n\ge4$ and let $\rho$ be a nontrivial ordinary irreducible
representation of $\operatorname{Sp}_{2n}(3)$ of degree $a$.
Every nonscalar finite-order operator $T$ normalizing its image
satisfies
\[
                              \nu(T)\ge\frac{13}{108}a.
                                                               \tag{HOS3.1}
\]
Both central quotients and all compatible outer automorphisms
are included.
\end{lemma}

\begin{proof}
The same complete universal-cover and automorphism inputs as in
D-HOS-09 apply, now with $q=3$: the universal cover is the
symplectic group and all automorphisms of the simple quotient are
inner-diagonal. There is no field or graph automorphism in this
range. The perfectness argument in (HOS5.15) therefore supplies a
natural $h\in\operatorname{GSp}_{2n}(3)$ implementing $T$ exactly
on the represented group; put $\mu=\mu(h)\in\{1,-1\}$.

Let $s$ be minus the identity on a nondegenerate natural plane and
the identity on its orthogonal complement. Choosing the line
radical from that plane, D-HOS-08 with $g=-I$ on its full
quotient gives
\[
               \nu(\rho(s))\ge\frac{1-3^{-(n-1)}}4a
                                      \ge\frac{13}{54}a.
                                                               \tag{HOS3.2}
\]
For any noncentral natural involution, choose the radical plane
in its smaller nonzero eigenspace. The larger eigenspace has
dimension $2m$ with $m\ge2$. The same argument gives
\[
                            \nu(\rho(s))\ge\frac29a.         \tag{HOS3.3}
\]

Suppose some $v$ has $b=B(v,hv)\ne0$. For any nonzero
$c\in\mathbf F_3$ the commutator
\[
 k=[h,t_v(c)]=t_{hv}(c/\mu)t_v(-c)
\]
acts as the identity outside $P=\langle v,hv\rangle$. Its
determinant on $P$ is one and its trace is
$2+c^2b^2/\mu=2+\mu^{-1}$. If $\mu=1$, this trace is zero
and $k^2=-I$ on $P$. If $\mu=-1$, its characteristic polynomial
on $P$ is $(X+1)^2$, and in characteristic three $k^3=-I$ on
$P$. Thus in either case a power is the rank-two involution
in (HOS3.2). The ordinary rank-commutator inequality (HOS5.17)
then gives $2\nu(T)\ge13a/54$.

Now suppose $B(v,hv)=0$ for all $v$. Polarization again gives
self-adjointness and $h^2=\mu I$. For $\mu=1$, Schur's lemma
makes $T$ a scalar multiple of a represented natural involution.
The central case would make $T$ scalar; otherwise (HOS3.3)
is stronger than the required bound.

It remains $\mu=-1$. For any $v\ne0$, put $r_1=v,r_2=hv$.
These are independent and span a totally isotropic plane. For
nonzero $c$, the same commutator is
\[
 k=I+N,\qquad
 N=-c\{r_1B(r_1,-)+r_2B(r_2,-)\},\quad
                 N^2=0,\quad\operatorname{rank}N=2.         \tag{HOS3.4}
\]
We need the equal coefficients in this expression, rather than
the absolute-trace bound for an arbitrary rank-two shear.

Complete the image plane to a nondegenerate natural four-space
$E_0$ as in D-HOS-09, and choose a paired $e,f$ in its
orthogonal complement. For the full line radical $Q_e$, the
induced quotient action is $g=k|_W$, of order three and moving
rank two. Choose symplectic coordinates $W=X\oplus Y$ with
$r_1,r_2$ the first two coordinate vectors in $Y$. Writing
$m=n-1$, this action is
\[
 g(x,y)=\bigl(x,y+c(x_1,x_2,0,\ldots,0)\bigr).
                                                               \tag{HOS3.5}
\]
In the explicit Schr\"odinger model (HOS.5), multiplication by
\[
 (\Omega f)(u)=
   \psi\bigl(-c(u_1^2+u_2^2)/2\bigr)f(u)
                                                               \tag{HOS3.6}
\]
implements (HOS3.5) exactly. Indeed conjugating $L_{x,y}$ by
the displayed multiplication operator gives $L_{x,y+cx'}$,
where $x'=(x_1,x_2,0,\ldots,0)$; the quadratic constant term
is exactly the required half-pairing term in $L_{x,y+cx'}$.

For either nontrivial central character, the values of the
quadratic exponent on its first two coordinates have
multiplicities $1,4,4$: the sum of two squares over
$\mathbf F_3$ is zero only at $(0,0)$, is one at four
vectors, and is minus one at four vectors. Multiplication by
a nonzero coefficient only interchanges the latter values.
The other $m-2$ coordinates are free. Thus every eigenspace
of this carrier has dimension at most $4/9$ of its degree.
As in D-HOS-08, tensoring with its arbitrary multiplicity
operator preserves this upper bound.

Among the nonzero linear characters of $Q_e$, all
multiplicities are equal under the full Levi action. The
proportion fixed by $g$ is less than $1/9$, and the other
orbits have length three. Hence their total support fraction
is at least $(2/3)(8/9)=16/27>5/9$. Combining the linear and
nontrivial-central-character summands, and then using
D-HOS-07, proves
\[
              \nu(\rho(k))\ge\frac59(a-\dim V^{Q_e})
                                     \ge\frac5{18}a.
                                                               \tag{HOS3.7}
\]
Equation (HOS5.17) gives $\nu(T)\ge5a/36>13a/108$.
This exhausts the self-adjoint case and proves the lemma.
\end{proof}

\begin{theorem}[D-HOS-12: selfdual mixed three-and-five cores]
\label{D-HOS-12}
In the imaginary-quadratic Weil-sixfold source problem, an inducing
image with an irreducibly acting normal central product of factors
of the following kinds cannot have period rank nine:
\begin{enumerate}
\item simple quotient $\operatorname{PSp}_{2n}(3)$, $n\ge4$,
with a nontrivial selfdual ordinary irreducible core representation;
\item simple quotient $\operatorname{PSp}_{2n}(5)$, $n\ge4$,
with any nontrivial ordinary irreducible core representation.
\end{enumerate}
Ranks and the two field sizes may be mixed; there is no transitivity
or original/inducing-field-equality assumption. Both polarization
determinants are included.
\end{theorem}

\begin{proof}
All core representations in the statement are selfdual; for the
second kind this was proved in D-HOS-10 from the published
anti-symplectic factorization. The elementary degree bound
D-HOS-02 gives degree at least $27$ for each factor of the first
kind and at least $250$ for each of the second. Thus the total
inducing degree $e$ is at least $27$.

The factorwise convolution and weighted tensor-cycle argument
in D-HOS-10 applies unchanged. Because
$13/108<31/250<1/4$, every nonscalar element of the full
compatible normalizer has projective support at least $13e/108$.
The unique core-invariant bilinear line is preserved by the
inducing image. Its similitude character is nontrivial because
the original imaginary quadratic character field is contained
in the inducing character field. The proof of these two
assertions is precisely the bilinear-form and character-induction
argument in D-HOS-10; it does not require that each inducing
character value already lie in the original quadratic field.

The actual intermediate Hurwitz factorization forces genus zero
and at least twelve genuine inertias, since $e\ge27$ and the
multiplicity variation has dimension six. At least two inertias
have nontrivial similitude and moving rank at least
$\lceil e/2\rceil$. Each remaining genuine inertia has moving
rank at least $\lceil13e/108\rceil$, including the full-rank
nonidentity scalar case. Consequently
\[
 2e+6\ge2\lceil e/2\rceil
                     +10\lceil13e/108\rceil.                \tag{HOS3.8}
\]
Dropping ceilings first gives $e\le324/11<30$. The elementary
degree lower bound leaves only $e=27,28,29$. For all three,
$\lceil13e/108\rceil=4$, so the right side of (HOS3.8) is at
least $68$, whereas the left side is at most $64$. This is a
contradiction. Nonselfdual higher $\mathbf F_3$ factors have
not been excluded by this argument.
\end{proof}
